\documentclass[11pt]{amsart}
\usepackage{calc,amssymb,amsmath,ulem,stmaryrd,fullpage}
\usepackage{enumerate}
\usepackage{alltt}
\RequirePackage[dvipsnames,usenames]{color}
\let\mffont=\sf
\normalem
%\input{kmacros3.sty}
\input{xy}
\xyoption{all}
\usepackage{hyperref}
\usepackage{graphicx}
\usepackage[all,cmtip]{xy}
\usepackage{verbatim}
\usepackage[left=0.95in,top=0.95in,right=0.95in,bottom=0.95in]{geometry}
\newcommand{\tauCohomology}{T}
\newcommand{\FCohomology}{S}


\theoremstyle{theorem}
%\newtheorem*{mainthm*}{Main Theorem}

\newcommand{\note}[1]{\marginpar{\sffamily\tiny #1}}

\def\todo#1{\textcolor{Mahogany}%
{\sffamily\footnotesize\newline{\color{Mahogany}\fbox{\parbox{\textwidth-15pt}{#1}}}\newline}}

\renewcommand{\O}{\mathcal O}
\newcommand{\ie}{{\itshape i.e.} }
\newcommand{\roundup}[1]{\lceil #1 \rceil}
\newcommand{\rounddown}[1]{\lfloor #1 \rfloor}
\renewcommand{\to}[1][]{\xrightarrow{\ #1\ }}
\newcommand{\onto}[1][]{\protect{\xrightarrow{\ #1\ }\hspace{-0.8em}\rightarrow}}
\newcommand{\into}[1][]{\lhook \joinrel \xrightarrow{\ #1\ }}
%\newcommand{DBDual}[1]{{\underline{\omega}_{#1}}}
\newcommand{\DBDual}[1]{{{\uuline \omega}{}^{\mydot}_{#1}}}
\newcommand{\Tt}{{\mathfrak{T}}}
%\newcommand{\bn}{{\mathfrak{n}} }
\newcommand{\sol}[1]{ \vskip 6pt{\color{blue} {\bf Solution:} #1} \vskip 6pt}

\newcommand{\bR}{{\mathbb{R}}}
\newcommand{\va}{{\vec{a}}}
\newcommand{\vb}{{\vec{b}}}
\newcommand{\vzero}{{\vec{0}}}
\newcommand{\vi}{{\vec{i}}}
\newcommand{\vj}{{\vec{j}}}
\newcommand{\vk}{{\vec{k}}}
\newcommand{\vh}{{\vec{h}}}
\newcommand{\vr}{{\vec{r}}}
\newcommand{\vu}{{\vec{u}}}
\newcommand{\vv}{{\vec{v}}}
\newcommand{\vw}{{\vec{w}}}
\newcommand{\vx}{{\vec{x}}}
\newcommand{\vy}{{\vec{y}}}
\newcommand{\vz}{{\vec{z}}}
\newcommand{\vT}{{\vec{T}}}
\newcommand{\vN}{{\vec{N}}}
\newcommand{\vF}{{\vec{F}}}
\newcommand{\vG}{{\vec{G}}}
\newcommand{\bF}{\mathbb{F}}
\newcommand{\bQ}{\mathbb{Q}}
\newcommand{\bZ}{\mathbb{Z}}
\newcommand{\bC}{\mathbb{C}}
\newcommand{\Gal}{\textnormal{Gal}}
\newcommand{\Aut}{\textnormal{Aut}}
\newcommand{\Emb}{\textnormal{Emb}}
\newcommand{\Tr}{\textnormal{Tr}}
\newcommand{\N}{\textnormal{N}}
\newcommand{\Hom}{\textnormal{Hom}}
\newcommand{\Spec}{\textnormal{Spec}}


\begin{document}
\title{HW \#6 -- Math 6320\\Spring 2015}
\author{Due: Thursday April 2nd}
\maketitle

\begin{enumerate}
\item  Suppose that $R$ is a Noetherian ring.  Is it true that there exists an integer $n_0 > 0$ such that
\begin{itemize}
\item[(a)]  every ideal $I \subseteq R$ is generated by at most $n_0$ elements?
\item[(b)]  every ascending chain of ideals $I_1 \subsetneq I_2 \subsetneq \ldots$ has length at most $n_0$?
\end{itemize}
Prove or give a counter example.
\item Let $A = k[x,y]$, $I = \langle x \rangle$ and $B = k$.  We have the natural projection $f : A \to A/I$ and the natural inclusion $g : k \hookrightarrow A/I \cong k[y]$.  Consider the ring $C = \{ (a,b) \in A \oplus B \; |\; f(a) = g(b) \}$ as in problem (8) from homework 5.  
    \begin{itemize}
    \item[(a)]  Show that projection the map from $i : C \to A$ is injective
    \item[(b)]  Show that $C$ is non-Noetherian even though it is a subring of a Noetherian ring.
    \item[(c)]  Describe the maximal ideal of $C$ in terms of the map $i^{\sharp}: \Spec A \to \Spec C$.  
    \end{itemize}
    \vskip 3pt
    \emph{Hint:} You may assume (8) from the previous homework.
\item Show that every finitely generated module over a Noetherian ring is finitely presented.
\item Show that every Artinian integral domain is a field.
\item An $R$-module $I$ is called \emph{injective} if for every \emph{injective} map of $R$-modules $f : A \to B$ and every $R$-module map $g : A \to I$, there exists a map $h : B \to I$ such that the following diagram commutes:
    \[
    \xymatrix{
    I \\
    A \ar[u]^g \ar[r]_f & B \ar@{.>}[ul]_h.
    }
    \]
    \begin{itemize}
    \item[(a)]  Show that $\bQ$ is an injective $\bZ$-module, as is $\bQ/\bZ$.
    \item[(b)]  Suppose that $0 \to I \to B \to C \to 0$ is exact, prove that
    \[
    0 \to \Hom_R(N, I) \to \Hom_R(N, B) \to \Hom_R(N, C) \to 0
    \]
    is exact for every $R$-module $N$ if $I$ is injective.
    \end{itemize}
\item Prove the 5-lemma.  In other words, suppose that
\[
\xymatrix{
0 \ar[r] & M_1 \ar[d]_{\alpha_1} \ar[r] &M_2 \ar[d]_{\alpha_2} \ar[r] &M_3 \ar[d]_{\alpha_3} \ar[r] &M_4 \ar[d]_{\alpha_4} \ar[r] &M_5 \ar[d]_{\alpha_5} \ar[r] & 0\\
0 \ar[r] & N_1 \ar[r] &N_2 \ar[r] &N_3 \ar[r] &N_4 \ar[r] &N_5 \ar[r] & 0
}
\]
is a commutative diagram with exact sequences as rows.
\begin{itemize}
\item[(a)]  If we suppose that $\alpha_2$ and $\alpha_4$ are surjective, and $\alpha_5$ is injective, show that $\alpha_3$ is surjective.
\item[(b)]  If we suppose that $\alpha_2$ and $\alpha_4$ are injective, and $\alpha_1$ is surjective, show that $\alpha_3$ is injective.
\item[(c)]  Show that if $\alpha_2, \alpha_4$ are isomorphisms, $\alpha_5$ is injective and $\alpha_1$ is surjective, that $\alpha_3$ is also an isomorphism.
\end{itemize}
\end{enumerate}

\end{document}

